
\subsection{Bayesian Inference}

In this section, we outline several important domains in which probabilistic modeling and statistical inference play a central role.  
So far, we have concentrated on analyzing abstract probability models, but in practice these models are useful only to the extent that they are relevant to, and informed by, real data.  
The field of inference and statistics is precisely concerned with using data generated by the real world to construct, select, and validate such models.  

Traditionally, statisticians would often be asked to analyze relatively small data sets, for example a limited number of medical records in order to evaluate whether a treatment is effective.  
In contrast, modern applications generate massive data sets and call for complex models involving a large number of parameters.  
At the same time, advances in computing power allow us to handle these large models, so that the range of applications and opportunities has expanded dramatically.

One classical setting involves the deliberate design of a data collection procedure, followed by a relatively simple prediction task.  
A basic example is polling, where we collect responses from a sample of individuals and use them to predict the outcome of an election.  
A closely related area is marketing and advertising, where the goal is to make predictions at the level of individual consumers rather than an entire population.  

A particularly prominent application in this context is that of recommendation systems.  
One collects data such as ratings that users assign to movies or other items, which can be organized in a large table with users indexing the rows and items indexing the columns.  
Using these partially observed data, one then aims to predict missing entries in the table, such as whether a given user would like a particular movie that they have not yet rated.

Finance is another major application domain.  
Financial markets are inherently uncertain, but they also produce extensive historical data sets.  
A central question is how to use these data to build models that can capture market behavior and support tasks such as risk assessment or pricing of financial instruments.  

In the natural sciences, there has been a revolution driven by the availability of large-scale data.  
In the life sciences, for instance, genomic data are analyzed in order to understand how combinations of genes relate to diseases or to map out complex biochemical pathways inside cells.  
An emerging frontier is neuroscience, where increasingly detailed measurements of brain activity raise the possibility of building models that describe how neural systems process information.  

Similar themes arise across many other scientific disciplines.  
In climate and environmental science, large and intricate models are calibrated using extensive observational data to understand and predict phenomena such as climate change.  
In physics, sophisticated inference techniques are used to detect rare events or objects, such as elusive particles or distant exoplanets, amidst large volumes of noisy measurements.  

Engineering applications also provide numerous examples.  
Engineers are often concerned with designing systems that perform reliably in uncertain or noisy environments, and probabilistic models guide both analysis and design.  
Signal processing is a prototypical area, where the goal is to recover an underlying signal from corrupted observations, such as reconstructing the content of a radio transmission received in the presence of noise.  

The range of application domains is enormous, and it continues to grow as new types of data become available.  
Despite the diversity of these settings, the fundamental probabilistic and statistical concepts that arise are remarkably similar.  
Our goal in the subsequent sections is to develop core methodologies that can be applied across this wide spectrum of problems.

\bigskip
\textsc{Illustrative example}

Consider the setting of movie recommendation systems.  
We can represent the data by a large table, where each row corresponds to a user and each column corresponds to a movie, and the entry in row \(i\) and column \(j\) is the rating that user \(i\) assigned to movie \(j\) whenever such a rating is available.  
The inference task is then to use the observed entries of this table to predict missing ratings, for example to decide whether user \(i\) is likely to enjoy movie \(j\) and should therefore be recommended that movie.


\subsubsection{Types of inference problems}

In the field of inference, it is useful to distinguish between different types of problems that arise.  
One broad distinction is between building models of a system and using a given model to make inferences about unknown quantities.  
In some situations, constructing or calibrating the model itself is the primary goal, while in others the model is taken as given and the focus is on unobserved variables that appear in that model.  

To illustrate this distinction, consider a simple communication setting.  
A transmitter sends a signal \(S\) through a medium, which attenuates the signal by a factor \(a\), and the signal is additionally corrupted by random noise \(W\).  
The receiver observes a signal \(X\) that can be described by the linear relation
\[
X = aS + W.
\]
Here, \(W\) is modeled as random noise with some known distribution, while \(a\) and \(S\) may be unknown.  

In one version of the problem, our goal is to characterize the medium.  
We transmit a known pilot signal \(S\), we observe the corresponding \(X\), and we use the relation \(X = aS + W\) to infer the value of the attenuation factor \(a\).  
In this case, we are effectively estimating a parameter of the system and thereby constructing or refining a model of the medium.  

In a different version of the problem, we assume that the medium has already been characterized.  
The attenuation factor \(a\) is known, but now the transmitted signal \(S\) is unknown to the receiver.  
Given the observation \(X\) and the model \(X = aS + W\), the task is to infer the value of \(S\).  

From a mathematical standpoint, these two versions have the same structural form.  
In both cases, we work with the same linear relation between \(X\), \(S\), and \(a\), but we treat different quantities as known or unknown.  
This means that similar methodologies can often be applied to both types of problems, even though their interpretations are quite different.

Another important distinction in inference is between hypothesis testing problems and estimation problems.  
In hypothesis testing, the unknown takes values in a finite or small discrete set, corresponding to a few candidate models of the world.  
The goal is to decide which of these candidate models is the correct one, based on the available data, while keeping the probability of making an incorrect decision small.  

A classical example is a radar detection problem.  
We receive a radar return signal and wish to decide between two hypotheses: either there is an airplane present in the range of the radar, or there is no airplane.  
The inference task is to choose between these two discrete alternatives, using the observed data and an appropriate decision rule.  

In estimation problems, the unknown is typically numerical and may vary over a continuum of possible values.  
We aim to produce an estimate that is close to the true value of the quantity of interest, according to some distance or loss function.  
Here, instead of choosing among a few discrete models, we select an estimate from an infinite set of possible values, such as the real line.  

These two types of problems, hypothesis testing and estimation, often require different solution strategies and performance criteria.  
In hypothesis testing, performance is usually quantified in terms of error probabilities associated with incorrect decisions between the competing hypotheses.  
In estimation, the emphasis is on the accuracy of the estimate, measured for example by mean squared error or another notion of distance from the true value.  

Despite these differences, both hypothesis testing and estimation are central components of statistical inference.  
Many practical problems involve elements of both, such as first deciding which model or scenario is likely to hold, and then estimating numerical parameters within the chosen model.  
In the following sections, we will develop methods tailored to each of these problem types and explore how they can be combined in more complex applications.

\bigskip
\textsc{Example: detection versus estimation}

Consider again the radar setting.  
In a pure detection problem, we use the observed signal to decide between the hypotheses \(\textit{airplane present}\) and \(\textit{airplane absent}\).  
In a related estimation problem, assuming that an airplane is present, we might further use the radar data to estimate a continuous quantity such as the distance or velocity of the airplane, aiming for an estimate that is close to the true value.


\subsubsection{The Bayesian inference framework}

In the Bayesian inference framework, the unknown quantity of interest is modeled as a random variable, denoted by \(\Theta\).  
Unlike other approaches in which \(\Theta\) is treated as an unknown but fixed constant, here \(\Theta\) has a probability distribution that captures our uncertainty about its value.  
This initial distribution is called the \textit{prior} distribution of \(\Theta\), and it encodes our beliefs about \(\Theta\) before any data are observed.  

We then obtain data in the form of an observation \(X\).  
The observation \(X\) is itself a random variable, and we model it probabilistically by specifying the conditional distribution of \(X\) given \(\Theta\).  
This conditional distribution describes how the data would be generated under each possible value of \(\Theta\), and it is sometimes referred to as the \textit{likelihood model}.  

Once we observe a specific numerical value \(x\) of \(X\), we can use Bayes' rule to update our prior beliefs.  
The result is the \textit{posterior} distribution of \(\Theta\) given the observation \(X = x\).  
In a discrete setting, this posterior is a conditional PMF, whereas in a continuous setting, it is a conditional PDF.  

The posterior distribution represents a complete Bayesian solution to the inference problem.  
It summarizes all the information about \(\Theta\) that is available after combining the prior with the observed data through the probabilistic model.  
In particular, it allows us to quantify the relative plausibility of different values of \(\Theta\) in light of the data.  

A natural question is how to choose the prior distribution.  
In some situations, symmetry or lack of prior information suggests that all values in a certain range should be regarded as equally likely, leading to a uniform prior over that range.  
In other cases, prior information may come from previous experiments or studies, and the new data can then be used to refine and update these earlier conclusions.  

The choice of prior can also be subjective, reflecting personal or expert beliefs about the parameter \(\Theta\).  
Regardless of its origin, the prior should respect any known constraints, such as the possible range of \(\Theta\).  
Values outside that range are assigned zero prior probability, ensuring that the posterior also assigns zero probability to impossible parameter values.  

To appreciate the idea of a posterior distribution, consider an election forecasting problem.  
Let \(\Theta\) denote the number of electoral votes that a candidate will receive.  
Instead of making a single yes/no prediction about the winner, a Bayesian analysis produces a full distribution over the possible values of \(\Theta\), providing a detailed picture of the uncertainty about the outcome.  

In such a setting, the results of a poll might be presented as a histogram or curve showing the probabilities of different possible vote totals.  
This graphical representation corresponds to the posterior distribution of \(\Theta\) given the polling data.  
It is much more informative than a single statement of which candidate is more likely to win, because it reveals how confident or uncertain the prediction really is.  

While the posterior distribution is a complete answer to the Bayesian inference problem, we often wish to summarize it by a single number called an \textit{estimate}.  
One common choice is the \textit{maximum a posteriori} (MAP) estimate.  
In the discrete case, the MAP estimate is the value of \(\Theta\) that has the largest posterior probability, and in the continuous case it is the point at which the posterior density is maximized.  

Another important summary is the \textit{conditional expectation} estimator.  
Here we report the posterior mean, that is, the conditional expectation \(\mathbb{E}[\Theta \mid X = x]\).  
This estimator is also known as the least mean squares (LMS) estimator, because it minimizes the mean squared error between the estimate and the true value of \(\Theta\).  

It is useful to distinguish between an \textit{estimator} and an \textit{estimate}.  
An estimator is a rule or function, typically denoted by \(g\), that maps observed data to a numerical value; in abstract terms, it is a random variable \(\hat{\Theta} = g(X)\).  
An estimate is the numerical value that results when we apply this rule to a particular observed value \(x\), namely \(\hat{\theta} = g(x)\).  

Because \(X\) is random, the estimator \(\hat{\Theta}\) is also a random variable and can be studied using probabilistic tools.  
Different estimators, such as the MAP estimator or the conditional expectation estimator, have different properties and performance measures.  
In subsequent discussions, we will analyze these properties and see how they guide the choice of an appropriate estimator for a given inference problem.

\bigskip
\textsc{Example: Bayesian estimation in an election}

Let \(\Theta\) be the (unknown) number of electoral votes that candidate A will receive in an election.  
Based on past elections and prior knowledge, we assign a prior distribution to \(\Theta\), reflecting which totals we consider more or less plausible before any new polling data are collected.  
After running a poll and observing data \(X = x\), we update our beliefs using Bayes' rule to obtain the posterior distribution of \(\Theta\), which can be visualized as a probability mass function over the possible vote totals and used to report quantities such as the MAP estimate or the posterior mean.


\subsubsection{Discrete parameter, discrete observation}

We now consider Bayesian inference in the case where both the unknown parameter \(\Theta\) and the observation \(X\) are discrete random variables.  
The parameter \(\Theta\) takes values in a finite or countable set, and we can interpret each possible value as a different hypothesis about the underlying model.  
Given a prior PMF for \(\Theta\) and a conditional PMF for \(X\) given \(\Theta\), Bayes' rule allows us to compute the posterior PMF of \(\Theta\) once an observation has been made.  

Suppose that \(\Theta\) can take values in the set \(\{1,2,3\}\).  
After observing a particular value \(X = x\), we apply Bayes' rule and obtain the conditional PMF \(p_{\Theta \mid X}(\theta \mid x)\).  
For concreteness, imagine that this posterior distribution assigns probabilities \(0.1\), \(0.6\), and \(0.3\) to the values \(\theta = 1\), \(\theta = 2\), and \(\theta = 3\), respectively.  

At this stage, the posterior distribution provides a complete probabilistic description of our uncertainty about \(\Theta\) given the observation.  
However, we may wish to summarize this information by reporting a single numerical estimate of \(\Theta\).  
One natural way to do this is to use the maximum a posteriori (MAP) rule, which selects the value of \(\theta\) that has the largest posterior probability.  

Formally, the MAP estimate \(\hat{\theta}_{\text{MAP}}\) is defined by
\[
\hat{\theta}_{\text{MAP}}
=
\arg\max_{\theta} \, p_{\Theta \mid X}(\theta \mid x).
\]
In the example above, the posterior probability is largest at \(\theta = 2\), so the MAP estimate would be \(\hat{\theta} = 2\).  
This estimate can be viewed as the most likely value of \(\Theta\) given the observed data.  

Another common way to construct an estimate is to use the least mean squares (LMS) rule.  
In this approach, we choose as our estimate the conditional expectation of \(\Theta\) given the observed value,
\[
\hat{\theta}_{\text{LMS}}
=
\mathbb{E}[\Theta \mid X = x]
=
\sum_{\theta} \theta \, p_{\Theta \mid X}(\theta \mid x).
\]
In the numerical example with probabilities \(0.1\), \(0.6\), and \(0.3\), this gives
\[
\hat{\theta}_{\text{LMS}}
=
1 \cdot 0.1 + 2 \cdot 0.6 + 3 \cdot 0.3
=
2.2.
\]

We now turn to the question of how to assess the quality of an estimate.  
When the values of \(\Theta\) are interpreted as distinct hypotheses, a natural performance criterion is the probability of error.  
Given that we have observed \(X = x\) and produced an estimate \(\hat{\theta}\), we define the conditional probability of error as
\[
\mathbb{P}\big(\hat{\Theta} \neq \Theta \mid X = x\big).
\]

Once the observation \(x\) has been fixed, the estimate \(\hat{\theta}\) is just a number, while \(\Theta\) remains random, described by the posterior distribution.  
The event of making an error is \(\{\Theta \neq \hat{\theta}\}\), and its conditional probability is
\[
\mathbb{P}\big(\Theta \neq \hat{\theta} \mid X = x\big)
=
1 - p_{\Theta \mid X}(\hat{\theta} \mid x).
\]
Equivalently, this is the sum of the posterior probabilities of all values of \(\theta\) different from \(\hat{\theta}\).  

In the numerical example, if we use the MAP estimate \(\hat{\theta} = 2\), the conditional probability of error becomes
\[
\mathbb{P}\big(\Theta \neq 2 \mid X = x\big)
=
0.1 + 0.3
=
0.4.
\]
If instead we had chosen \(\hat{\theta} = 3\), the conditional probability of error would be
\[
\mathbb{P}\big(\Theta \neq 3 \mid X = x\big)
=
0.1 + 0.6
=
0.7,
\]
which is larger.  

More generally, for any fixed observation \(x\), choosing \(\hat{\theta}\) to maximize the posterior probability \(p_{\Theta \mid X}(\hat{\theta} \mid x)\) is equivalent to minimizing the conditional probability of error.  
Indeed, minimizing
\(\mathbb{P}(\Theta \neq \hat{\theta} \mid X = x)\) is the same as maximizing \(\mathbb{P}(\Theta = \hat{\theta} \mid X = x)\).  
Thus, the MAP rule is optimal in the sense that it produces the smallest possible conditional probability of error for each given value of the observation.  

Often, we are also interested in the overall probability of error, not conditioned on a specific observation.  
This overall probability is
\[
\mathbb{P}\big(\hat{\Theta} \neq \Theta\big),
\]
which averages the conditional error probabilities over all possible observations.  
Using the law of total probability and conditioning on the observation, we can write
\[
\mathbb{P}\big(\hat{\Theta} \neq \Theta\big)
=
\sum_{x}
\mathbb{P}\big(\hat{\Theta} \neq \Theta \mid X = x\big)
\mathbb{P}(X = x).
\]

There is also an alternative expression obtained by conditioning on the value of \(\Theta\) instead of \(X\).  
In that case, the overall probability of error becomes
\[
\mathbb{P}\big(\hat{\Theta} \neq \Theta\big)
=
\sum_{\theta}
\mathbb{P}\big(\hat{\Theta} \neq \theta \mid \Theta = \theta\big)
\mathbb{P}(\Theta = \theta).
\]
Either of these formulas can be used for calculations, and the choice between them depends on which conditioning leads to simpler expressions in a given problem.  

The optimality of the MAP rule for the overall probability of error follows from its optimality at the conditional level.  
For each value of \(x\), the MAP estimate minimizes \(\mathbb{P}(\hat{\Theta} \neq \Theta \mid X = x)\).  
Since the overall probability of error is a weighted sum of these conditional probabilities, it must also be smallest when we use the MAP rule.  

In summary, when \(\Theta\) is discrete and interpreted as a set of alternative hypotheses, the MAP rule is the preferred estimator when the objective is to minimize the probability of making an incorrect decision.  
The LMS estimator, on the other hand, is more appropriate when the parameter is numerical and performance is measured in terms of mean squared error.  
Both estimators arise naturally from the posterior distribution, which encapsulates all the information provided by the data and the prior model.

\bigskip
\textsc{Example: three-hypothesis decision}

Suppose that \(\Theta\) takes values in \(\{1,2,3\}\) and that, after observing \(X = x\), the posterior PMF is given by
\[
p_{\Theta \mid X}(\theta \mid x)
=
\begin{cases}
0.1, & \theta = 1,\\
0.6, & \theta = 2,\\
0.3, & \theta = 3.
\end{cases}
\]
The MAP estimate is \(\hat{\theta}_{\text{MAP}} = 2\), because this value has the largest posterior probability.  
The corresponding conditional probability of error is
\[
\mathbb{P}\big(\Theta \neq 2 \mid X = x\big)
=
0.1 + 0.3
=
0.4,
\]
which is smaller than the error probability that would result from choosing either \(\hat{\theta} = 1\) or \(\hat{\theta} = 3\).


\newpage
\subsubsection{Discrete parameter, continuous observation}

We now consider Bayesian inference in the case where the parameter \(\Theta\) is discrete, while the observation \(X\) is continuous.  
As before, \(\Theta\) takes values in a finite or countable set, and we can interpret each possible value as a different hypothesis about the underlying model.  
The key difference is that the distribution of \(X\) is now described by probability density functions rather than probability mass functions.  

The prior information about \(\Theta\) is captured by a PMF \(p_{\Theta}(\theta)\).  
For the observation \(X\), we specify, for each possible value of \(\Theta\), a conditional probability density function \(f_{X \mid \Theta}(x \mid \theta)\).  
These conditional densities describe how the continuous observation is generated under each hypothesis.  

Bayes' rule still provides the link between the prior and the posterior, with the only change being that PDFs now appear in place of PMFs for the observation.  
For a given observation \(X = x\), the posterior PMF of \(\Theta\) is given by
\[
p_{\Theta \mid X}(\theta \mid x)
=
\frac{f_{X \mid \Theta}(x \mid \theta)\,p_{\Theta}(\theta)}
{\sum_{\theta'} f_{X \mid \Theta}(x \mid \theta')\,p_{\Theta}(\theta')},
\]
where the denominator is the unconditional PDF of \(X\) evaluated at \(x\).  

A standard example arises in communication systems.  
Suppose that \(\Theta\) can take one of three possible values, corresponding to three different signals that might be transmitted.  
The transmitted signal is corrupted by additive noise, and the receiver observes a random variable
\[
X = \Theta + W,
\]
where \(W\) represents noise.  

A typical assumption is that the noise \(W\) has mean zero, some known variance, and is independent of \(\Theta\).  
If we know the PDF \(f_{W}(w)\), then conditioning on \(\Theta = \theta\) yields
\[
X \mid (\Theta = \theta) = \theta + W.
\]
Adding the constant \(\theta\) simply shifts the distribution of \(W\), so the conditional density of \(X\) given \(\Theta = \theta\) is
\[
f_{X \mid \Theta}(x \mid \theta)
=
f_{W}(x - \theta).
\]

Using these conditional densities together with the prior \(p_{\Theta}(\theta)\), we can apply Bayes' rule and compute the posterior probabilities \(p_{\Theta \mid X}(\theta \mid x)\).  
Once this posterior PMF has been obtained, the situation is analogous to the discrete–discrete case.  
In particular, we can again apply the maximum a posteriori rule to construct an estimate of \(\Theta\).  

The MAP estimate \(\hat{\theta}_{\text{MAP}}\) is defined, for the observed value \(x\), by
\[
\hat{\theta}_{\text{MAP}}
=
\arg\max_{\theta} \, p_{\Theta \mid X}(\theta \mid x).
\]
That is, we choose the value of \(\theta\) that has the largest posterior probability given the continuous observation.  
This value is interpreted as the most likely hypothesis in light of the data.  

Once an estimate has been chosen, we may again be interested in the probability of error.  
For a fixed observation \(x\), the conditional probability of error is
\[
\mathbb{P}\big(\hat{\Theta} \neq \Theta \mid X = x\big)
=
1 - p_{\Theta \mid X}(\hat{\theta} \mid x),
\]
since an error occurs whenever \(\Theta\) differs from the chosen estimate \(\hat{\theta}\).  
As in the purely discrete case, minimizing this conditional probability of error is equivalent to maximizing the posterior probability of the chosen hypothesis.  

Thus, for each fixed \(x\), the MAP rule produces the smallest possible conditional probability of error.  
The fact that \(X\) is continuous does not change this conclusion.  
The posterior PMF \(p_{\Theta \mid X}(\theta \mid x)\) plays exactly the same conceptual role as in the discrete–discrete setting.  

We can also consider the overall, or unconditional, probability of error,
\[
\mathbb{P}\big(\hat{\Theta} \neq \Theta\big),
\]
which averages the conditional probabilities of error over all possible observations and parameters.  
One way to express this quantity is to condition on the observation \(X\) and integrate over all possible values of \(x\).  

In this form, the law of total probability becomes
\[
\mathbb{P}\big(\hat{\Theta} \neq \Theta\big)
=
\int
\mathbb{P}\big(\hat{\Theta} \neq \Theta \mid X = x\big)
f_{X}(x)\,dx,
\]
where \(f_{X}(x)\) is the unconditional PDF of \(X\).  
This expression is the continuous analogue of the sum that appears in the discrete case.  

An alternative expression is obtained by conditioning on the parameter \(\Theta\).  
We can write
\[
\mathbb{P}\big(\hat{\Theta} \neq \Theta\big)
=
\sum_{\theta}
\mathbb{P}\big(\hat{\Theta} \neq \theta \mid \Theta = \theta\big)
p_{\Theta}(\theta).
\]
In practical calculations, it may be more convenient to use either the integral representation or this summation representation, depending on which is easier to evaluate for the given model.  

The optimality of the MAP rule for the overall probability of error follows the same reasoning as in the discrete–discrete case.  
For each fixed observation \(x\), the MAP rule minimizes \(\mathbb{P}(\hat{\Theta} \neq \Theta \mid X = x)\).  
Since the overall probability of error is obtained by integrating these conditional probabilities with respect to the distribution of \(X\), the MAP rule also minimizes the unconditional probability of error.  

Therefore, when \(\Theta\) represents a finite set of hypotheses and the performance criterion is the probability of making an incorrect decision, the MAP rule is optimal whether the observation \(X\) is discrete or continuous.  
The role of the continuous observation is to determine, through the likelihood functions and Bayes' rule, the posterior probabilities on the discrete set of hypotheses.  
Once those posterior probabilities are known, the decision problem and its solution proceed exactly as in the purely discrete setting.

\bigskip
\textsc{Example: three-level signal in noise}

Suppose that a transmitter sends a signal that can take three possible levels \(\{-1,0,1\}\), and let \(\Theta\) denote the transmitted level.  
The received signal is
\[
X = \Theta + W,
\]
where \(W\) is zero-mean Gaussian noise with variance \(\sigma^{2}\) and is independent of \(\Theta\).  
Given an observed value \(x\), the conditional density \(f_{X \mid \Theta}(x \mid \theta)\) is a Gaussian PDF centered at \(\theta\), and Bayes' rule combines these densities with the prior probabilities \(p_{\Theta}(\theta)\) to produce the posterior probabilities \(p_{\Theta \mid X}(\theta \mid x)\); the MAP estimate then selects the value of \(\theta \in \{-1,0,1\}\) with the largest posterior probability, yielding the decision rule that minimizes the probability of symbol detection error.


\subsubsection{Continuous parameter, continuous observation}

We now turn to the case where both the parameter \(\Theta\) and the observation \(X\) are continuous random variables.  
As before, the Bayesian framework specifies a prior distribution for \(\Theta\) and a conditional distribution for \(X\) given \(\Theta\).  
Bayes' rule then provides the posterior distribution of \(\Theta\) given the observed value of \(X\).  

In this setting, all distributions are described by probability density functions.  
The prior information about the parameter is encoded in a prior density \(f_{\Theta}(\theta)\).  
The observation model is given by the conditional density \(f_{X \mid \Theta}(x \mid \theta)\), which describes how \(X\) would be distributed if the parameter took the value \(\theta\).  

Once we observe a particular numerical value \(x\), Bayes' rule takes the form
\[
f_{\Theta \mid X}(\theta \mid x)
=
\frac{f_{X \mid \Theta}(x \mid \theta)\,f_{\Theta}(\theta)}
{\int f_{X \mid \Theta}(x \mid \theta')\,f_{\Theta}(\theta')\,d\theta'},
\]
which defines the posterior density of \(\Theta\) given \(X = x\).  
This posterior density summarizes our updated knowledge about \(\Theta\) after taking into account both the prior and the observed data.  

A particularly important and widely used class of continuous Bayesian models is the family of linear normal models.  
In such models, the variables of interest are combined through linear relationships, and all of them are modeled as Gaussian random variables.  
Because Gaussian distributions are closed under linear transformations, these models often allow closed-form expressions for the posterior distribution.  

As a simple example, consider a continuous-valued signal \(\Theta\) that we wish to estimate.  
We observe a noisy version of this signal,
\[
X = \Theta + W,
\]
where \(W\) is additive noise that is independent of \(\Theta\) and has a known distribution, often taken to be normal with mean zero and some variance.  
Given the observation \(X = x\), our task is to infer the underlying signal value \(\Theta\).  

More elaborate versions of this problem involve vector-valued parameters and multiple observations.  
In such cases, \(\Theta\) may be a vector with several components, and we may obtain a collection of measurements \(X\) that depend linearly on \(\Theta\) plus noise.  
These multivariate linear normal models form an important part of the study of Bayesian inference and will be examined in more detail elsewhere.  

Another example of a continuous model arises when estimating the parameter of a uniform distribution.  
Suppose that \(X\) is uniformly distributed over an interval whose endpoints depend on an unknown parameter \(\Theta\).  
The parameter \(\Theta\) is then given a prior distribution, and the goal is to use one or more observations of \(X\) to infer the true value of \(\Theta\).  

In the continuous–continuous setting, we again wish to construct estimators of \(\Theta\) based on the observation \(X\).  
Two principal Bayesian estimators appear here, just as in the discrete case.  
The first is the maximum a posteriori estimator, and the second is the least mean squares estimator.  

Given the posterior density \(f_{\Theta \mid X}(\theta \mid x)\), the maximum a posteriori (MAP) estimator chooses the value of \(\theta\) at which this density is largest.  
For a specific observed value \(x\), the MAP estimate is defined by
\[
\hat{\theta}_{\text{MAP}}
=
\arg\max_{\theta} \, f_{\Theta \mid X}(\theta \mid x).
\]
This provides a single value that is regarded as the most plausible parameter value in light of the data and the prior.  

The least mean squares (LMS) estimator is based on the conditional expectation of \(\Theta\) given \(X\).  
For each observed value \(x\), the LMS estimate is
\[
\hat{\theta}_{\text{LMS}}
=
\mathbb{E}[\Theta \mid X = x]
=
\int \theta \, f_{\Theta \mid X}(\theta \mid x)\,d\theta.
\]
This estimator is also known as the conditional expectation estimator or the posterior mean.  

To evaluate and compare estimators in this setting, a natural performance measure is the squared distance between the estimate and the true parameter value.  
For a given estimator \(\hat{\Theta} = g(X)\) and a given observation \(X = x\), the conditional mean squared error is defined as
\[
\mathbb{E}\big[(\hat{\Theta} - \Theta)^{2} \mid X = x\big].
\]
This quantity measures, on average, how far the estimator is from the true parameter, conditioned on the observed data.  

We can also consider the unconditional, or overall, mean squared error of the estimator.  
This is obtained by averaging over both the randomness in \(\Theta\) and the randomness in \(X\),
\[
\mathbb{E}\big[(\hat{\Theta} - \Theta)^{2}\big],
\]
and it provides a single number summarizing the overall accuracy of the estimator.  
Both conditional and unconditional mean squared error play an important role in analyzing the performance of estimators.  

The LMS estimator has a particularly important optimality property with respect to mean squared error.  
Among all possible estimators that are functions of \(X\), the conditional expectation \(\mathbb{E}[\Theta \mid X]\) achieves the smallest mean squared error.  
This fact justifies the interpretation of the LMS estimator as the best, in the mean squared error sense, within the Bayesian framework.  

In practice, the choice between the MAP estimator and the LMS estimator depends on the goals of the application.  
If we are interested in the most probable parameter value, the MAP estimate is a natural choice.  
If instead we wish to minimize mean squared error, the LMS estimator is the appropriate one.  

In summary, the continuous–continuous Bayesian setting extends the same ideas we have seen in discrete models to the realm of probability densities.  
We specify a prior, a likelihood model, and then apply Bayes' rule to obtain the posterior distribution.  
From this posterior, we can construct estimators such as the MAP and LMS estimators and evaluate their performance using criteria such as mean squared error.

\bigskip
\textsc{Example: Gaussian signal in Gaussian noise}

Let \(\Theta\) be a Gaussian random variable with mean \(\mu_{\Theta}\) and variance \(\sigma_{\Theta}^{2}\).  
Suppose we observe
\[
X = \Theta + W,
\]
where \(W\) is independent Gaussian noise with mean \(0\) and variance \(\sigma_{W}^{2}\).  
In this linear normal model, the posterior distribution of \(\Theta\) given \(X = x\) is also Gaussian, and both the MAP and LMS estimators coincide, yielding a linear function of \(x\) that minimizes the mean squared error between the estimate and the true value of \(\Theta\).


\subsubsection{Inferring the unknown bias of a coin and the Beta distribution}

We consider a coin with an unknown bias \(\Theta\), where \(\Theta\) takes values in the interval \([0,1]\).  
Consistent with the Bayesian viewpoint, we treat \(\Theta\) as a random variable and assign to it a prior distribution.  
We then flip the coin \(n\) times independently and record the number of heads obtained.  

Let \(K\) denote the number of heads in these \(n\) tosses.  
Given a particular value \(\Theta = \theta\), the random variable \(K\) has a binomial distribution with parameters \(n\) and \(\theta\).  
Our goal is to use the observed value of \(K\) to infer information about the unknown parameter \(\Theta\).  

We begin with the case where the prior distribution of \(\Theta\) is uniform on the unit interval.  
This prior can be interpreted as expressing a state of ignorance, in the sense that all values of \(\theta \in [0,1]\) are considered equally plausible a priori.  
The corresponding prior density is
\[
f_{\Theta}(\theta) = 1,
\qquad
0 \leq \theta \leq 1.
\]

Once we observe a specific value \(K = k\), Bayes' rule in the mixed continuous–discrete setting takes the form
\[
f_{\Theta \mid K}(\theta \mid k)
=
\frac{p_{K \mid \Theta}(k \mid \theta)\,f_{\Theta}(\theta)}
{\int_{0}^{1} p_{K \mid \Theta}(k \mid \theta')\,f_{\Theta}(\theta')\,d\theta'},
\]
for \(0 \leq \theta \leq 1\).  
Here, \(f_{\Theta \mid K}(\theta \mid k)\) is the posterior density of \(\Theta\) given that \(K = k\).  

Given \(\Theta = \theta\), the number of heads \(K\) follows a binomial distribution, so
\[
p_{K \mid \Theta}(k \mid \theta)
=
\binom{n}{k} \theta^{k} (1 - \theta)^{n-k},
\qquad
k = 0,1,\ldots,n.
\]
Substituting this into Bayes' rule and using the uniform prior, we obtain
\[
f_{\Theta \mid K}(\theta \mid k)
=
\frac{\binom{n}{k} \theta^{k} (1 - \theta)^{n-k} \cdot 1}
{\int_{0}^{1} \binom{n}{k} \theta'^{k} (1 - \theta')^{n-k}\,d\theta'},
\qquad
0 \leq \theta \leq 1.
\]

The binomial coefficient \(\binom{n}{k}\) does not depend on \(\theta\) and can be absorbed into the normalizing constant in the denominator.  
Thus, up to a constant factor, the dependence on \(\theta\) in the posterior density is
\[
f_{\Theta \mid K}(\theta \mid k)
\propto
\theta^{k} (1 - \theta)^{n-k},
\qquad
0 \leq \theta \leq 1.
\]
Outside the interval \([0,1]\), both the prior and the posterior densities are equal to zero.  

Densities of the form
\[
f(\theta)
\propto
\theta^{\alpha - 1} (1 - \theta)^{\beta - 1},
\qquad
0 \leq \theta \leq 1,
\]
are known as \textit{Beta} distributions with parameters \(\alpha\) and \(\beta\).  
Comparing this with the expression above, we see that the posterior distribution of \(\Theta\) given \(K = k\) is a Beta distribution with parameters \(\alpha = k + 1\) and \(\beta = n - k + 1\).  

Next, we consider the more general case in which the prior distribution of \(\Theta\) is itself a Beta distribution.  
Specifically, suppose that
\[
f_{\Theta}(\theta)
\propto
\theta^{\alpha - 1} (1 - \theta)^{\beta - 1},
\qquad
0 \leq \theta \leq 1,
\]
for some \(\alpha > 0\) and \(\beta > 0\).  
This prior allows us to encode prior beliefs that favor certain regions of the interval \([0,1]\).  

Using the same likelihood as before, the joint dependence on \(\theta\) in the numerator of Bayes' rule is
\[
p_{K \mid \Theta}(k \mid \theta)\,f_{\Theta}(\theta)
\propto
\Big[\theta^{k} (1 - \theta)^{n-k}\Big]
\Big[\theta^{\alpha - 1} (1 - \theta)^{\beta - 1}\Big].
\]
Multiplying these terms, we obtain
\[
p_{K \mid \Theta}(k \mid \theta)\,f_{\Theta}(\theta)
\propto
\theta^{\alpha + k - 1} (1 - \theta)^{\beta + n - k - 1},
\qquad
0 \leq \theta \leq 1.
\]

After normalization, the posterior density takes the form
\[
f_{\Theta \mid K}(\theta \mid k)
\propto
\theta^{\alpha + k - 1} (1 - \theta)^{\beta + n - k - 1},
\qquad
0 \leq \theta \leq 1,
\]
and is equal to zero outside this interval.  
This shows that the posterior distribution of \(\Theta\) is again a Beta distribution, now with updated parameters \(\alpha' = \alpha + k\) and \(\beta' = \beta + n - k\).  

We have thus discovered an important property of the Beta family.  
If the prior distribution of \(\Theta\) is Beta with parameters \(\alpha\) and \(\beta\), and the data consist of \(n\) Bernoulli trials with \(k\) heads, then the posterior distribution of \(\Theta\) is also Beta, with parameters \(\alpha + k\) and \(\beta + n - k\).  
In other words, the Beta distribution is \textit{conjugate} to the Bernoulli and binomial likelihood in the Bayesian framework.  

This conjugacy has several practical advantages.  
It leads to simple update rules for the parameters \(\alpha\) and \(\beta\) as more data are observed.  
In particular, each observed head increases the first parameter \(\alpha\) by one, and each observed tail increases the second parameter \(\beta\) by one.  

Because of this structure, it is convenient to update the posterior distribution recursively as new coin tosses are recorded.  
Starting from an initial Beta prior, we can process each observation in turn, updating the parameters after each toss instead of recomputing the posterior from scratch.  
This makes the Beta–binomial model a powerful and flexible tool for modeling and learning about unknown probabilities in Bernoulli-type experiments.

\bigskip
\textsc{Example: updating a Beta prior}

Suppose that the prior distribution of the coin bias \(\Theta\) is Beta with parameters \(\alpha = 2\) and \(\beta = 3\), reflecting prior information that favors values of \(\theta\) near \(0.4\).  
We then flip the coin \(n = 5\) times and observe \(k = 4\) heads.  
Using the update rule, the posterior distribution of \(\Theta\) is Beta with parameters \(\alpha' = 2 + 4 = 6\) and \(\beta' = 3 + 5 - 4 = 4\), which concentrates more mass near higher values of \(\theta\) and reflects the evidence from the observed data.


\subsubsection{Inferring the unknown bias of a coin: point estimates}

We continue with the problem of estimating the unknown bias \(\Theta\) of a coin, where \(\Theta \in [0,1]\).  
We assume a uniform prior on \([0,1]\) and observe the number of heads \(K\) in \(n\) independent tosses of the coin.  
The posterior density of \(\Theta\) given \(K = k\) is a Beta distribution of the form
\[
f_{\Theta \mid K}(\theta \mid k)
=
d(n,k)\,\theta^{k}(1-\theta)^{n-k},
\qquad
0 \leq \theta \leq 1,
\]
where \(d(n,k)\) is a normalizing constant that depends on \(n\) and \(k\) but not on \(\theta\).  

We first derive the maximum a posteriori (MAP) estimate of \(\Theta\).  
By definition, the MAP estimate is the value of \(\theta\) at which the posterior density \(f_{\Theta \mid K}(\theta \mid k)\) attains its maximum.  
Since \(d(n,k)\) does not depend on \(\theta\), maximizing the posterior is equivalent to maximizing the function
\[
g(\theta)
=
\theta^{k}(1-\theta)^{n-k},
\qquad
0 < \theta < 1.
\]

It is more convenient to maximize the logarithm of this function.  
Define
\[
\ell(\theta)
=
\log g(\theta)
=
k \log \theta
+
(n-k)\log(1-\theta).
\]
To find the maximizer, we differentiate \(\ell(\theta)\) with respect to \(\theta\) and set the derivative equal to zero.  

The derivative is
\[
\ell'(\theta)
=
\frac{k}{\theta}
-
\frac{n-k}{1-\theta}.
\]
Setting \(\ell'(\theta) = 0\) gives
\[
\frac{k}{\theta}
=
\frac{n-k}{1-\theta},
\]
which implies
\[
k(1-\theta)
=
\theta(n-k).
\]
Solving for \(\theta\), we obtain
\[
k - k\theta
=
n\theta - k\theta
\quad\Rightarrow\quad
k
=
n\theta
\quad\Rightarrow\quad
\theta
=
\frac{k}{n}.
\]

Thus, for the observed value \(K = k\), the MAP estimate of the bias is
\[
\hat{\theta}_{\text{MAP}}
=
\frac{k}{n}.
\]
This expression is quite natural: it coincides with the empirical frequency of heads in the \(n\) tosses.  
The corresponding estimator, viewed as a function of the random variable \(K\), is
\[
\hat{\Theta}_{\text{MAP}}
=
\frac{K}{n},
\]
which is a random variable whose realized value is the MAP estimate once \(K\) is observed.  

We now turn to the least mean squares (LMS) estimator, which is the conditional expectation of \(\Theta\) given \(K = k\).  
By definition,
\[
\hat{\theta}_{\text{LMS}}
=
\mathbb{E}[\Theta \mid K = k]
=
\int_{0}^{1} \theta \, f_{\Theta \mid K}(\theta \mid k)\,d\theta.
\]
Substituting the expression for the posterior density, we obtain
\[
\hat{\theta}_{\text{LMS}}
=
\int_{0}^{1} \theta \, d(n,k)\,\theta^{k}(1-\theta)^{n-k}\,d\theta
=
d(n,k)
\int_{0}^{1} \theta^{k+1}(1-\theta)^{n-k}\,d\theta.
\]

To proceed, we need an explicit expression for the normalizing constant \(d(n,k)\).  
This constant is determined by the requirement that \(f_{\Theta \mid K}(\theta \mid k)\) integrates to 1 over the interval \([0,1]\).  
Thus,
\[
1
=
\int_{0}^{1} f_{\Theta \mid K}(\theta \mid k)\,d\theta
=
d(n,k)
\int_{0}^{1} \theta^{k}(1-\theta)^{n-k}\,d\theta,
\]
so
\[
d(n,k)
=
\frac{1}
{\int_{0}^{1} \theta^{k}(1-\theta)^{n-k}\,d\theta}.
\]

We will make use of the following integral identity.  
For non-negative integers \(\alpha\) and \(\beta\),
\[
\int_{0}^{1} \theta^{\alpha}(1-\theta)^{\beta}\,d\theta
=
\frac{\alpha!\,\beta!}{(\alpha + \beta + 1)!}.
\]
Applying this formula with \(\alpha = k\) and \(\beta = n-k\), we find
\[
\int_{0}^{1} \theta^{k}(1-\theta)^{n-k}\,d\theta
=
\frac{k!\,(n-k)!}{(n+1)!},
\]
and therefore
\[
d(n,k)
=
\frac{(n+1)!}{k!\,(n-k)!}.
\]

We now evaluate the integral appearing in the expression for \(\hat{\theta}_{\text{LMS}}\).  
Using the same identity with \(\alpha = k+1\) and \(\beta = n-k\), we obtain
\[
\int_{0}^{1} \theta^{k+1}(1-\theta)^{n-k}\,d\theta
=
\frac{(k+1)!\,(n-k)!}{(n+2)!}.
\]
Substituting into the expression for \(\hat{\theta}_{\text{LMS}}\), we get
\[
\hat{\theta}_{\text{LMS}}
=
d(n,k)
\,
\frac{(k+1)!\,(n-k)!}{(n+2)!}
=
\frac{(n+1)!}{k!\,(n-k)!}
\,
\frac{(k+1)!\,(n-k)!}{(n+2)!}.
\]

We simplify this expression step by step.  
First, we cancel the factors \((n-k)!\) in the numerator and denominator, obtaining
\[
\hat{\theta}_{\text{LMS}}
=
\frac{(n+1)!}{k!}
\,
\frac{(k+1)!}{(n+2)!}.
\]
Next, we use the identities
\[
(k+1)! = (k+1)\,k!,
\qquad
(n+2)! = (n+2)\,(n+1)!,
\]
which lead to
\[
\hat{\theta}_{\text{LMS}}
=
\frac{(n+1)!}{k!}
\,
\frac{(k+1)\,k!}{(n+2)\,(n+1)!}
=
\frac{k+1}{n+2}.
\]

Thus, the LMS estimate of \(\Theta\) given \(K = k\) is
\[
\hat{\theta}_{\text{LMS}}
=
\frac{k+1}{n+2}.
\]
The corresponding estimator, as a function of the random variable \(K\), is
\[
\hat{\Theta}_{\text{LMS}}
=
\frac{K+1}{n+2}.
\]

We can now compare the MAP and LMS estimates.  
The MAP estimate for the observed value \(K = k\) is
\[
\hat{\theta}_{\text{MAP}} = \frac{k}{n},
\]
while the LMS estimate is
\[
\hat{\theta}_{\text{LMS}} = \frac{k+1}{n+2}.
\]
These two expressions are similar but not identical, which reflects the fact that the mode and the mean of a Beta distribution are generally different.  

However, when the number of tosses \(n\) is large, the two estimates are close to each other.  
Indeed, for large \(n\), the difference between \(k/n\) and \((k+1)/(n+2)\) becomes negligible.  
In this sense, as the amount of data increases, the MAP and LMS estimators for the coin bias nearly coincide.  

\bigskip
\textsc{Example: numerical comparison of MAP and LMS estimates}

Suppose we flip a coin \(n = 8\) times and observe \(k = 6\) heads.  
Under a uniform prior, the MAP estimate of the bias is
\[
\hat{\theta}_{\text{MAP}}
=
\frac{6}{8}
=
0.75,
\]
while the LMS estimate is
\[
\hat{\theta}_{\text{LMS}}
=
\frac{6+1}{8+2}
=
\frac{7}{10}
=
0.7.
\]
In this moderate-sample regime, the two estimates are close but not identical, with the LMS estimate being slightly more conservative and closer to the center of the interval.



\subsubsection{Recognizing normal PDFs}

Recall that a normal random variable \(X \sim \mathcal{N}(\mu,\sigma^{2})\) has a PDF of the form
\[
f_{X}(x)
=
\frac{1}{\sqrt{2\pi\sigma^{2}}}
\exp\Big(-\frac{(x-\mu)^{2}}{2\sigma^{2}}\Big).
\]
Sometimes, we are given a density that already looks like a normal density, and we want to recognize its mean and variance by matching it to this standard form.

Suppose we are told that a random variable \(X\) has PDF
\[
f_{X}(x)
=
c \exp\Big(-8(x-3)^{2}\Big),
\]
for some normalizing constant \(c\).  
To identify the mean and variance, we compare the exponent with the canonical normal form.  
We can write
\[
-8(x-3)^{2}
=
-\frac{(x-3)^{2}}{2\sigma^{2}},
\]
which implies
\[
\frac{1}{2\sigma^{2}}
=
8
\quad\Rightarrow\quad
\sigma^{2}
=
\frac{1}{16}.
\]
The term \(x-3\) in the exponent shows that the distribution is centered at \(3\), so \(\mu = 3\).  
Finally, to find \(c\), we note that the PDF must integrate to \(1\) and that the normalizing constant for a normal distribution is \(1/(\sqrt{2\pi\sigma^{2}})\).  
Since \(\sigma^{2} = 1/16\), we have \(\sigma = 1/4\), and hence
\[
c
=
\frac{1}{\sqrt{2\pi\sigma^{2}}}
=
\frac{1}{\sigma\sqrt{2\pi}}
=
\frac{1}{\tfrac{1}{4}\sqrt{2\pi}}
=
\frac{4}{\sqrt{2\pi}}.
\]

\bigskip
\textsc{Recognizing a general quadratic exponent}

Consider now a more general PDF of the form
\[
f_{X}(x)
=
c \exp\big(-\alpha x^{2} - \beta x - \gamma\big),
\]
where \(\alpha > 0\) and \(c\) is a normalizing constant.  
We will show that this is also a normal PDF and identify its mean and variance.

The key step is to rewrite the quadratic term in the exponent by \textit{completing the square}.  
We start from
\[
\alpha x^{2} + \beta x + \gamma.
\]
Factor out \(\alpha\) from the first two terms:
\[
\alpha x^{2} + \beta x + \gamma
=
\alpha\Big(x^{2} + \frac{\beta}{\alpha}x\Big) + \gamma.
\]
We now complete the square inside the parentheses.  
Observe that
\[
x^{2} + \frac{\beta}{\alpha}x
=
\Big(x + \frac{\beta}{2\alpha}\Big)^{2}
-
\Big(\frac{\beta}{2\alpha}\Big)^{2}.
\]
Substituting this identity, we obtain
\[
\alpha x^{2} + \beta x + \gamma
=
\alpha\Big(x + \frac{\beta}{2\alpha}\Big)^{2}
-
\alpha\Big(\frac{\beta}{2\alpha}\Big)^{2}
+
\gamma
=
\alpha\Big(x + \frac{\beta}{2\alpha}\Big)^{2}
-
\frac{\beta^{2}}{4\alpha}
+
\gamma.
\]
Therefore, the exponent in the PDF can be rewritten as
\[
-\alpha x^{2} - \beta x - \gamma
=
-\alpha\Big(x + \frac{\beta}{2\alpha}\Big)^{2}
+
\frac{\beta^{2}}{4\alpha}
-
\gamma.
\]
This leads to the expression
\[
f_{X}(x)
=
c \exp\Big(-\alpha\Big(x + \frac{\beta}{2\alpha}\Big)^{2}\Big)
\exp\Big(\frac{\beta^{2}}{4\alpha} - \gamma\Big).
\]
The second exponential factor does not depend on \(x\), so it can be absorbed into the constant by defining
\[
c'
=
c \exp\Big(\frac{\beta^{2}}{4\alpha} - \gamma\Big).
\]
Thus, the PDF becomes
\[
f_{X}(x)
=
c'
\exp\Big(-\alpha\Big(x + \frac{\beta}{2\alpha}\Big)^{2}\Big).
\]

We can now match this with the standard normal form.  
We want to write the exponent as
\[
-\frac{(x-\mu)^{2}}{2\sigma^{2}}.
\]
Comparing this expression with
\[
-\alpha\Big(x + \frac{\beta}{2\alpha}\Big)^{2},
\]
we see that
\[
x - \mu
=
x + \frac{\beta}{2\alpha},
\]
so that
\[
\mu
=
-\frac{\beta}{2\alpha}.
\]
For the variance, we equate the coefficients of the squared term in the exponent:
\[
\alpha
=
\frac{1}{2\sigma^{2}}
\quad\Rightarrow\quad
\sigma^{2}
=
\frac{1}{2\alpha}.
\]
Therefore, \(X\) is normally distributed with mean \(\mu = -\beta/(2\alpha)\) and variance \(\sigma^{2} = 1/(2\alpha)\).  
The normalizing constant \(c'\) must then equal \(1/\sqrt{2\pi\sigma^{2}}\).  
From this, if desired, one can solve for the original constant \(c\).

\bigskip
\textsc{Using the peak to identify the mean}

Once we know that a density is normal, there is a faster way to identify its mean.  
The mean of a normal random variable coincides with the location of the peak of its PDF.  
Thus, rather than completing the square in full, we can instead find the point \(x\) where the exponent is maximized (or, equivalently, where the quadratic inside the minus sign is minimized).

For the PDF
\[
f_{X}(x)
=
c \exp\big(-\alpha x^{2} - \beta x - \gamma\big),
\]
the exponent is
\[
-\alpha x^{2} - \beta x - \gamma.
\]
Maximizing this expression is equivalent to minimizing the quadratic
\[
\alpha x^{2} + \beta x + \gamma.
\]
To find the minimizing point, we differentiate with respect to \(x\) and set the derivative equal to zero:
\[
\frac{d}{dx}\big(\alpha x^{2} + \beta x + \gamma\big)
=
2\alpha x + \beta
=
0.
\]
Solving for \(x\) gives
\[
x
=
-\frac{\beta}{2\alpha}.
\]
This is the location of the minimum of the quadratic, hence the location of the maximum of the exponent and the peak of the PDF.  
Since the distribution is normal, this point is exactly the mean:
\[
\mu
=
-\frac{\beta}{2\alpha}.
\]
The variance can still be recovered by matching the quadratic coefficient \(\alpha\) to the standard form, yielding again
\[
\sigma^{2}
=
\frac{1}{2\alpha}.
\]


\subsubsection{Normal unknown and additive noise}

We consider a simple linear model with additive noise.  
There is an unknown parameter, modeled as a random variable \(\Theta\), that we wish to estimate.  
What we observe is a noisy version of \(\Theta\), given by
\[
X
=
\Theta + W,
\]
where \(W\) is an additive noise term.

We assume that \(\Theta\) and \(W\) are independent standard normal random variables, that is,
\[
\Theta \sim \mathcal{N}(0,1),
\qquad
W \sim \mathcal{N}(0,1),
\qquad
\Theta \perp W.
\]
According to the Bayesian approach, inference about \(\Theta\) is based on the posterior distribution of \(\Theta\) given the observation \(X = x\).  
To obtain this posterior, we first write down the prior and the likelihood.

The prior density of \(\Theta\) is the standard normal density
\[
f_{\Theta}(\theta)
=
\frac{1}{\sqrt{2\pi}}
\exp\Big(-\frac{\theta^{2}}{2}\Big).
\]
Next, we determine the conditional distribution of \(X\) given \(\Theta\).  
Condition on the event \(\{\Theta = \theta\}\).  
In this conditional universe, the observation satisfies
\[
X
\mid (\Theta = \theta)
=
\theta + W.
\]
Since \(W\) is independent of \(\Theta\), conditioning on \(\Theta = \theta\) does not change the distribution of \(W\).  
Thus \(W\) remains standard normal, and \(X\) is the sum of a constant \(\theta\) and a standard normal random variable.  
This implies
\[
X \mid (\Theta = \theta)
\sim
\mathcal{N}(\theta,1).
\]

The corresponding conditional density of \(X\) given \(\Theta = \theta\) is
\[
f_{X\mid\Theta}(x\mid\theta)
=
\frac{1}{\sqrt{2\pi}}
\exp\Big(-\frac{(x-\theta)^{2}}{2}\Big).
\]
Bayes' rule gives the posterior density of \(\Theta\) given \(X=x\) as
\[
f_{\Theta\mid X}(\theta\mid x)
=
\frac{f_{X\mid\Theta}(x\mid\theta)\,f_{\Theta}(\theta)}{f_{X}(x)},
\]
where \(f_{X}(x)\) is the marginal density of \(X\).  
For purposes of identifying the posterior form, we can ignore the denominator \(f_{X}(x)\), which is just a normalizing constant depending on \(x\) but not on \(\theta\).  
Thus, up to a proportionality constant,
\[
f_{\Theta\mid X}(\theta\mid x)
\propto
f_{X\mid\Theta}(x\mid\theta)\,f_{\Theta}(\theta).
\]

Substituting the expressions for the prior and the likelihood, we obtain
\[
f_{\Theta\mid X}(\theta\mid x)
\propto
\exp\Big(-\frac{(x-\theta)^{2}}{2}\Big)
\exp\Big(-\frac{\theta^{2}}{2}\Big).
\]
Combining the exponents, we get
\[
f_{\Theta\mid X}(\theta\mid x)
\propto
\exp\Big(-\frac{(x-\theta)^{2} + \theta^{2}}{2}\Big).
\]
We now expand the quadratic in \(\theta\):
\[
(x-\theta)^{2} + \theta^{2}
=
(x^{2} - 2x\theta + \theta^{2}) + \theta^{2}
=
x^{2} - 2x\theta + 2\theta^{2}.
\]
Thus
\[
f_{\Theta\mid X}(\theta\mid x)
\propto
\exp\Big(-\frac{x^{2} - 2x\theta + 2\theta^{2}}{2}\Big).
\]
We rearrange the exponent by separating the terms that depend on \(\theta\) from the ones that depend only on \(x\):
\[
-\frac{x^{2} - 2x\theta + 2\theta^{2}}{2}
=
-\frac{x^{2}}{2}
+
x\theta
-
\theta^{2}.
\]
Hence
\[
f_{\Theta\mid X}(\theta\mid x)
\propto
\exp\Big(-\frac{x^{2}}{2}\Big)
\exp\big(x\theta - \theta^{2}\big).
\]
The factor \(\exp(-x^{2}/2)\) does not involve \(\theta\) and can be absorbed into the normalizing constant.  
Therefore, as a function of \(\theta\),
\[
f_{\Theta\mid X}(\theta\mid x)
\propto
\exp\big(x\theta - \theta^{2}\big).
\]

We rewrite the exponent \(x\theta - \theta^{2}\) by completing the square in \(\theta\).  
Observe that
\[
x\theta - \theta^{2}
=
-\theta^{2} + x\theta
=
-\Big(\theta^{2} - x\theta\Big)
=
-\Big[\theta^{2} - x\theta + \Big(\frac{x}{2}\Big)^{2} - \Big(\frac{x}{2}\Big)^{2}\Big]
=
-\Big(\theta - \frac{x}{2}\Big)^{2}
+
\Big(\frac{x}{2}\Big)^{2}.
\]
Substituting into the exponent gives
\[
f_{\Theta\mid X}(\theta\mid x)
\propto
\exp\Big(-\Big(\theta - \frac{x}{2}\Big)^{2} + \Big(\frac{x}{2}\Big)^{2}\Big)
=
\exp\Big(\Big(\frac{x}{2}\Big)^{2}\Big)
\exp\Big(-\Big(\theta - \frac{x}{2}\Big)^{2}\Big).
\]
Again, the factor \(\exp\big((x/2)^{2}\big)\) is a constant with respect to \(\theta\) and can be absorbed into the normalization.  
Thus the posterior density has the form
\[
f_{\Theta\mid X}(\theta\mid x)
\propto
\exp\Big(-\Big(\theta - \frac{x}{2}\Big)^{2}\Big).
\]
We recognize this as the kernel of a normal density with mean \(x/2\) and variance \(1/2\).  
Indeed, a normal density with mean \(m\) and variance \(\sigma^{2}\) has exponent \(-(\theta-m)^{2}/(2\sigma^{2})\).  
In our case, the exponent is \(-(\theta - x/2)^{2}\), which corresponds to
\[
\frac{1}{2\sigma^{2}}
=
1
\quad\Rightarrow\quad
\sigma^{2}
=
\frac{1}{2}.
\]
Therefore, the posterior distribution is
\[
\Theta \mid (X = x)
\sim
\mathcal{N}\Big(\frac{x}{2}, \frac{1}{2}\Big).
\]

Since the posterior is normal, the conditional expectation and the mode coincide.  
Thus, the least mean squares estimator (posterior mean) and the maximum a posteriori (MAP) estimator are the same:
\[
\mathbb{E}[\Theta \mid X = x]
=
\frac{x}{2}.
\]
Equivalently, we can obtain this value by locating the peak of the posterior density.  
Because the posterior kernel is
\[
\exp\big(x\theta - \theta^{2}\big),
\]
we can find the peak by minimizing the quadratic term inside the negative exponent.  
Ignoring constants that do not depend on \(\theta\), the relevant function is
\[
q(\theta)
=
\theta^{2} - x\theta.
\]
To minimize \(q(\theta)\), we take the derivative with respect to \(\theta\) and set it equal to zero:
\[
\frac{d}{d\theta}q(\theta)
=
2\theta - x
=
0.
\]
Solving for \(\theta\) yields
\[
\theta
=
\frac{x}{2}.
\]
This is the location of the maximum of the posterior density and therefore the posterior mean and MAP estimate.

From the perspective of estimators (random variables), we define the estimator \(\hat{\Theta}(X)\) as
\[
\hat{\Theta}(X)
=
\mathbb{E}[\Theta \mid X]
=
\frac{X}{2}.
\]
When we observe a particular numerical value \(X = x\), the corresponding estimate is the number \(x/2\).  
This estimator balances prior information (that \(\Theta\) is centered at \(0\)) with the observed data \(X\), producing an estimate that lies halfway between the prior mean \(0\) and the observation \(x\).

\bigskip
\textsc{Generalizations}

The same structure arises in more general settings where \(\Theta\) and \(W\) are independent normal random variables with arbitrary means and variances.  
In such cases, one can show, by repeating the same type of calculation, that the posterior distribution of \(\Theta\) given \(X\) is still normal.  
Because the posterior is normal, its mean coincides with its mode, so the least mean squares estimator and the MAP estimator are again identical.  
Moreover, the resulting estimator of \(\Theta\) turns out to be a linear function of the observation \(X\), reflecting a general feature of linear models with Gaussian priors and Gaussian noise.



\subsubsection{The case of multiple observations}

We consider a model in which an unknown parameter \(\Theta\) is observed through multiple noisy measurements.  
For \(i = 1,\dots,n\), each observation satisfies
\[
X_i
=
\Theta + W_i,
\]
where \(W_i\) is an additive noise term.

We assume that \(\Theta\) is a normal random variable with mean \(x_0\) and variance \(\sigma_{0}^{2}\):
\[
\Theta \sim \mathcal{N}(x_0,\sigma_{0}^{2}).
\]
For each \(i\), the noise term is also normal with mean \(0\) and variance \(\sigma_{i}^{2}\):
\[
W_i \sim \mathcal{N}(0,\sigma_{i}^{2}),
\]
and we assume that \(\Theta,W_1,\dots,W_n\) are all mutually independent.  
Our goal is to estimate \(\Theta\) on the basis of the observations \(X_1,\dots,X_n\).

In the Bayesian framework, estimation is based on the posterior distribution of \(\Theta\) given the observations.  
Let \(\mathbf{X} = (X_1,\dots,X_n)\) and \(\mathbf{x} = (x_1,\dots,x_n)\).  
Bayes' rule for continuous random variables gives
\[
f_{\Theta\mid \mathbf{X}}(\theta\mid \mathbf{x})
=
\frac{f_{\mathbf{X}\mid\Theta}(\mathbf{x}\mid\theta)\,f_{\Theta}(\theta)}{f_{\mathbf{X}}(\mathbf{x})},
\]
where \(f_{\Theta}\) is the prior density of \(\Theta\), \(f_{\mathbf{X}\mid\Theta}\) is the conditional joint density of the measurements given \(\Theta\), and \(f_{\mathbf{X}}\) is the marginal density of \(\mathbf{X}\).

We start by determining the conditional distribution of a single observation \(X_i\) given \(\Theta\).  
Condition on \(\{\Theta = \theta\}\).  
In this conditional universe, we have
\[
X_i \mid (\Theta = \theta)
=
\theta + W_i.
\]
Since \(W_i\) is independent of \(\Theta\), its distribution is unaffected by conditioning and remains \(\mathcal{N}(0,\sigma_{i}^{2})\).  
Thus
\[
X_i \mid (\Theta = \theta)
\sim
\mathcal{N}(\theta,\sigma_{i}^{2}).
\]
The conditional density of \(X_i\) given \(\Theta = \theta\) is therefore
\[
f_{X_i\mid\Theta}(x_i\mid\theta)
=
\frac{1}{\sqrt{2\pi\sigma_{i}^{2}}}
\exp\Big(-\frac{(x_i-\theta)^{2}}{2\sigma_{i}^{2}}\Big).
\]

We now turn to the conditional joint density \(f_{\mathbf{X}\mid\Theta}(\mathbf{x}\mid\theta)\).  
Given \(\Theta = \theta\), the measurement equations are
\[
X_i
=
\theta + W_i,
\quad
i=1,\dots,n.
\]
The noise variables \(W_1,\dots,W_n\) are independent, and conditioning on \(\Theta = \theta\) does not change this independence.  
Since each \(X_i\) is obtained from \(W_i\) by adding the constant \(\theta\), the random variables \(X_1,\dots,X_n\) are also independent in the conditional universe.  
Therefore, the conditional joint density factors as
\[
f_{\mathbf{X}\mid\Theta}(\mathbf{x}\mid\theta)
=
\prod_{i=1}^{n} f_{X_i\mid\Theta}(x_i\mid\theta)
=
\prod_{i=1}^{n}
\frac{1}{\sqrt{2\pi\sigma_{i}^{2}}}
\exp\Big(-\frac{(x_i-\theta)^{2}}{2\sigma_{i}^{2}}\Big).
\]

The prior density of \(\Theta\) is
\[
f_{\Theta}(\theta)
=
\frac{1}{\sqrt{2\pi\sigma_{0}^{2}}}
\exp\Big(-\frac{(\theta - x_0)^{2}}{2\sigma_{0}^{2}}\Big).
\]
Plugging these expressions into Bayes' rule, we obtain, up to a normalizing constant,
\[
f_{\Theta\mid \mathbf{X}}(\theta\mid \mathbf{x})
\propto
f_{\mathbf{X}\mid\Theta}(\mathbf{x}\mid\theta)\,f_{\Theta}(\theta)
=
\Bigg[
\prod_{i=1}^{n}
\frac{1}{\sqrt{2\pi\sigma_{i}^{2}}}
\exp\Big(-\frac{(x_i-\theta)^{2}}{2\sigma_{i}^{2}}\Big)
\Bigg]
\frac{1}{\sqrt{2\pi\sigma_{0}^{2}}}
\exp\Big(-\frac{(\theta - x_0)^{2}}{2\sigma_{0}^{2}}\Big).
\]
The prefactors involving \(1/\sqrt{2\pi\sigma_{i}^{2}}\) and \(1/\sqrt{2\pi\sigma_{0}^{2}}\) do not depend on \(\theta\) and can be absorbed into a constant.  
Thus
\[
f_{\Theta\mid \mathbf{X}}(\theta\mid \mathbf{x})
\propto
\exp\Bigg(
-\frac{(\theta - x_0)^{2}}{2\sigma_{0}^{2}}
-
\sum_{i=1}^{n}\frac{(x_i-\theta)^{2}}{2\sigma_{i}^{2}}
\Bigg).
\]

To identify the posterior distribution, we focus on the exponent as a function of \(\theta\).  
Let
\[
Q(\theta)
=
\frac{(\theta - x_0)^{2}}{2\sigma_{0}^{2}}
+
\sum_{i=1}^{n}\frac{(x_i-\theta)^{2}}{2\sigma_{i}^{2}}.
\]
Then
\[
f_{\Theta\mid \mathbf{X}}(\theta\mid \mathbf{x})
\propto
\exp\big(-Q(\theta)\big).
\]
We expand \(Q(\theta)\) and collect terms that are quadratic, linear, and constant in \(\theta\).  
First, for the prior term,
\[
\frac{(\theta - x_0)^{2}}{2\sigma_{0}^{2}}
=
\frac{1}{2\sigma_{0}^{2}}\big(\theta^{2} - 2x_0\theta + x_0^{2}\big)
=
\frac{\theta^{2}}{2\sigma_{0}^{2}}
-
\frac{x_0}{\sigma_{0}^{2}}\theta
+
\frac{x_0^{2}}{2\sigma_{0}^{2}}.
\]
Similarly, each measurement term expands as
\[
\frac{(x_i-\theta)^{2}}{2\sigma_{i}^{2}}
=
\frac{1}{2\sigma_{i}^{2}}\big(x_i^{2} - 2x_i\theta + \theta^{2}\big)
=
\frac{\theta^{2}}{2\sigma_{i}^{2}}
-
\frac{x_i}{\sigma_{i}^{2}}\theta
+
\frac{x_i^{2}}{2\sigma_{i}^{2}}.
\]
Summing over \(i\), we obtain
\[
\sum_{i=1}^{n}\frac{(x_i-\theta)^{2}}{2\sigma_{i}^{2}}
=
\theta^{2}\sum_{i=1}^{n}\frac{1}{2\sigma_{i}^{2}}
-
\theta\sum_{i=1}^{n}\frac{x_i}{\sigma_{i}^{2}}
+
\sum_{i=1}^{n}\frac{x_i^{2}}{2\sigma_{i}^{2}}.
\]

Combining the prior term and the measurement terms, the quadratic function \(Q(\theta)\) takes the form
\[
Q(\theta)
=
\frac{1}{2}\Bigg(
\Big(\frac{1}{\sigma_{0}^{2}} + \sum_{i=1}^{n}\frac{1}{\sigma_{i}^{2}}\Big)\theta^{2}
-
2\Big(\frac{x_0}{\sigma_{0}^{2}} + \sum_{i=1}^{n}\frac{x_i}{\sigma_{i}^{2}}\Big)\theta
\Big)
+
\text{(constant in \(\theta\))}.
\]
The constant term does not affect the shape of the posterior density as a function of \(\theta\), so we omit it.  
Thus, up to an additive constant,
\[
Q(\theta)
=
\frac{1}{2}A\theta^{2} - B\theta,
\]
where
\[
A
=
\frac{1}{\sigma_{0}^{2}} + \sum_{i=1}^{n}\frac{1}{\sigma_{i}^{2}},
\qquad
B
=
\frac{x_0}{\sigma_{0}^{2}} + \sum_{i=1}^{n}\frac{x_i}{\sigma_{i}^{2}}.
\]
Hence,
\[
f_{\Theta\mid \mathbf{X}}(\theta\mid \mathbf{x})
\propto
\exp\Big(-\frac{1}{2}A\theta^{2} + B\theta\Big).
\]

We now complete the square in \(\theta\).  
Consider the exponent
\[
-\frac{1}{2}A\theta^{2} + B\theta
=
-\frac{A}{2}\theta^{2} + B\theta.
\]
Factor out \(-\frac{A}{2}\):
\[
-\frac{A}{2}\theta^{2} + B\theta
=
-\frac{A}{2}\Big(\theta^{2} - \frac{2B}{A}\theta\Big).
\]
We write the quadratic inside the parentheses as a complete square:
\[
\theta^{2} - \frac{2B}{A}\theta
=
\Big(\theta - \frac{B}{A}\Big)^{2}
-
\Big(\frac{B}{A}\Big)^{2}.
\]
Substituting,
\[
-\frac{A}{2}\theta^{2} + B\theta
=
-\frac{A}{2}\Big(\Big(\theta - \frac{B}{A}\Big)^{2} - \Big(\frac{B}{A}\Big)^{2}\Big)
=
-\frac{A}{2}\Big(\theta - \frac{B}{A}\Big)^{2}
+
\frac{A}{2}\Big(\frac{B}{A}\Big)^{2}.
\]
Therefore,
\[
f_{\Theta\mid \mathbf{X}}(\theta\mid \mathbf{x})
\propto
\exp\Big(-\frac{A}{2}\Big(\theta - \frac{B}{A}\Big)^{2}\Big)
\exp\Big(\frac{A}{2}\Big(\frac{B}{A}\Big)^{2}\Big).
\]
The second exponential factor does not depend on \(\theta\) and can be absorbed into the normalizing constant.  
Thus the posterior density has the form
\[
f_{\Theta\mid \mathbf{X}}(\theta\mid \mathbf{x})
\propto
\exp\Big(-\frac{A}{2}\Big(\theta - \frac{B}{A}\Big)^{2}\Big).
\]
Comparing with the standard normal kernel
\[
\exp\Big(-\frac{(\theta - m)^{2}}{2\sigma^{2}}\Big),
\]
we identify
\[
\frac{1}{2\sigma^{2}}
=
\frac{A}{2}
\quad\Rightarrow\quad
\sigma^{2}
=
\frac{1}{A},
\]
and
\[
m
=
\frac{B}{A}.
\]
Therefore, the posterior distribution is normal:
\[
\Theta \mid \mathbf{X} = \mathbf{x}
\sim
\mathcal{N}\Big(\frac{B}{A}, \frac{1}{A}\Big),
\]
that is,
\[
\Theta \mid \mathbf{X} = \mathbf{x}
\sim
\mathcal{N}\Bigg(
\frac{\displaystyle \frac{x_0}{\sigma_{0}^{2}} + \sum_{i=1}^{n}\frac{x_i}{\sigma_{i}^{2}}}{\displaystyle \frac{1}{\sigma_{0}^{2}} + \sum_{i=1}^{n}\frac{1}{\sigma_{i}^{2}}},
\;
\frac{1}{\displaystyle \frac{1}{\sigma_{0}^{2}} + \sum_{i=1}^{n}\frac{1}{\sigma_{i}^{2}}}
\Bigg).
\]

Because the posterior is normal, the conditional expectation and the mode coincide.  
Thus, the least mean squares estimator (posterior mean) and the maximum a posteriori (MAP) estimator are identical.  
The posterior mean is
\[
\mathbb{E}[\Theta \mid \mathbf{X} = \mathbf{x}]
=
\frac{\displaystyle \frac{x_0}{\sigma_{0}^{2}} + \sum_{i=1}^{n}\frac{x_i}{\sigma_{i}^{2}}}{\displaystyle \frac{1}{\sigma_{0}^{2}} + \sum_{i=1}^{n}\frac{1}{\sigma_{i}^{2}}}.
\]
Interpreted as an estimator (a random variable), we write
\[
\hat{\Theta}(\mathbf{X})
=
\mathbb{E}[\Theta \mid \mathbf{X}]
=
\frac{\displaystyle \frac{x_0}{\sigma_{0}^{2}} + \sum_{i=1}^{n}\frac{X_i}{\sigma_{i}^{2}}}{\displaystyle \frac{1}{\sigma_{0}^{2}} + \sum_{i=1}^{n}\frac{1}{\sigma_{i}^{2}}}.
\]

This expression has a natural interpretation as a weighted average.  
Each term \(X_i\) is multiplied by the weight \(1/\sigma_{i}^{2}\), and the prior mean \(x_0\) is multiplied by the weight \(1/\sigma_{0}^{2}\).  
The denominator is the sum of all these weights, so we can write
\[
\hat{\Theta}(\mathbf{X})
=
\sum_{i=0}^{n} w_i Z_i,
\]
where
\[
Z_0 = x_0,\quad Z_i = X_i\ \text{for } i\ge 1,
\]
and
\[
w_i
=
\frac{1/\sigma_{i}^{2}}{\displaystyle \sum_{j=0}^{n}1/\sigma_{j}^{2}},
\quad
i = 0,\dots,n,
\]
with the convention that \(\sigma_{0}^{2}\) is the prior variance.  
The prior mean \(x_0\) thus acts as an additional observation, and all pieces of information are combined in a symmetric and linear fashion.

The dependence of the weights on the variances matches our intuition.  
If a particular variance \(\sigma_{i}^{2}\) is large, the corresponding observation \(X_i\) is very noisy and should receive a small weight.  
Conversely, observations with small noise variance contribute more to the final estimate.  
In summary, the Bayesian estimator in this normal model is a linear function of the measurements and the prior mean, and the posterior distribution of \(\Theta\) remains normal.



\subsubsection{The mean squared error}

We continue with the model in which an unknown parameter \(\Theta\) is observed through multiple noisy measurements.  
For \(i = 1,\dots,n\),
\[
X_i
=
\Theta + W_i,
\]
where \(\Theta \sim \mathcal{N}(x_0,\sigma_{0}^{2})\), \(W_i \sim \mathcal{N}(0,\sigma_{i}^{2})\), and \(\Theta,W_1,\dots,W_n\) are independent.  
We have seen that the posterior distribution of \(\Theta\) given the observations \(\mathbf{X} = (X_1,\dots,X_n)\) is normal with density
\[
f_{\Theta\mid \mathbf{X}}(\theta\mid \mathbf{x})
\propto
\exp\Bigg(
-\frac{(\theta - x_0)^{2}}{2\sigma_{0}^{2}}
-
\sum_{i=1}^{n}\frac{(x_i-\theta)^{2}}{2\sigma_{i}^{2}}
\Bigg).
\]

The posterior mean, which is also the maximum a posteriori (MAP) estimate, is
\[
\hat{\theta}
=
\mathbb{E}[\Theta \mid \mathbf{X} = \mathbf{x}]
=
\frac{\displaystyle \frac{x_0}{\sigma_{0}^{2}} + \sum_{i=1}^{n}\frac{x_i}{\sigma_{i}^{2}}}{\displaystyle \frac{1}{\sigma_{0}^{2}} + \sum_{i=1}^{n}\frac{1}{\sigma_{i}^{2}}}.
\]
As an estimator (a function of the random vector \(\mathbf{X}\)), we write
\[
\hat{\Theta}(\mathbf{X})
=
\frac{\displaystyle \frac{x_0}{\sigma_{0}^{2}} + \sum_{i=1}^{n}\frac{X_i}{\sigma_{i}^{2}}}{\displaystyle \frac{1}{\sigma_{0}^{2}} + \sum_{i=1}^{n}\frac{1}{\sigma_{i}^{2}}}.
\]

We now assess the performance of this estimator in terms of the mean squared error (MSE).  
Given that we have observed \(\mathbf{X} = \mathbf{x}\), the estimator \(\hat{\Theta}(\mathbf{X})\) takes the numerical value \(\hat{\theta}\), and the remaining uncertainty about \(\Theta\) is described by the posterior distribution.  
The conditional mean squared error is defined as
\[
\mathbb{E}\big[(\Theta - \hat{\Theta}(\mathbf{X}))^{2} \,\big|\, \mathbf{X} = \mathbf{x}\big].
\]
In the conditional universe where \(\mathbf{X} = \mathbf{x}\) is fixed, \(\hat{\Theta}(\mathbf{X}) = \hat{\theta}\) is a constant.  
Moreover, \(\hat{\theta}\) is the mean of the posterior distribution of \(\Theta\) given \(\mathbf{X} = \mathbf{x}\).  
Thus we are taking the expectation of the squared distance from the mean:
\[
\mathbb{E}\big[(\Theta - \hat{\Theta}(\mathbf{X}))^{2} \,\big|\, \mathbf{X} = \mathbf{x}\big]
=
\mathbb{E}\big[(\Theta - \hat{\theta})^{2} \,\big|\, \mathbf{X} = \mathbf{x}\big],
\]
which is exactly the variance of the posterior distribution:
\[
\mathbb{E}\big[(\Theta - \hat{\Theta}(\mathbf{X}))^{2} \,\big|\, \mathbf{X} = \mathbf{x}\big]
=
\operatorname{Var}(\Theta \mid \mathbf{X} = \mathbf{x}).
\]
Therefore, the conditional mean squared error is the posterior variance of \(\Theta\).

We now compute this posterior variance.  
Recall that the exponent in the posterior density can be written as a quadratic function of \(\theta\):
\[
-\frac{(\theta - x_0)^{2}}{2\sigma_{0}^{2}}
-
\sum_{i=1}^{n}\frac{(x_i-\theta)^{2}}{2\sigma_{i}^{2}}
=
-\frac{1}{2}A\theta^{2} + B\theta + \text{(constant)},
\]
where
\[
A
=
\frac{1}{\sigma_{0}^{2}} + \sum_{i=1}^{n}\frac{1}{\sigma_{i}^{2}}.
\]
From the previous derivation, we know that
\[
\Theta \mid \mathbf{X} = \mathbf{x}
\sim
\mathcal{N}\Big(\frac{B}{A},\frac{1}{A}\Big).
\]
Thus the posterior variance is
\[
\operatorname{Var}(\Theta \mid \mathbf{X} = \mathbf{x})
=
\frac{1}{A}
=
\frac{1}{\displaystyle \frac{1}{\sigma_{0}^{2}} + \sum_{i=1}^{n}\frac{1}{\sigma_{i}^{2}}}.
\]
This is the conditional mean squared error of the estimator \(\hat{\Theta}(\mathbf{X})\) given the specific observations \(\mathbf{X} = \mathbf{x}\):
\[
\mathbb{E}\big[(\Theta - \hat{\Theta}(\mathbf{X}))^{2} \,\big|\, \mathbf{X} = \mathbf{x}\big]
=
\frac{1}{\displaystyle \frac{1}{\sigma_{0}^{2}} + \sum_{i=1}^{n}\frac{1}{\sigma_{i}^{2}}}.
\]
A remarkable property of this model is that this quantity does not depend on the particular observed values \(x_1,\dots,x_n\); it depends only on the variances.  
Thus, regardless of what data we observe, the remaining uncertainty about \(\Theta\), measured by the posterior variance, is always the same.

We now consider the unconditional mean squared error, that is, the MSE before performing any measurements.  
This is
\[
\mathbb{E}\big[(\Theta - \hat{\Theta}(\mathbf{X}))^{2}\big].
\]
Using the law of total expectation,
\[
\mathbb{E}\big[(\Theta - \hat{\Theta}(\mathbf{X}))^{2}\big]
=
\mathbb{E}\Big[\mathbb{E}\big[(\Theta - \hat{\Theta}(\mathbf{X}))^{2} \,\big|\, \mathbf{X}\big]\Big].
\]
For each realization \(\mathbf{X} = \mathbf{x}\), the inner expectation is the conditional MSE, which we have found to be equal to the constant
\[
\frac{1}{\displaystyle \frac{1}{\sigma_{0}^{2}} + \sum_{i=1}^{n}\frac{1}{\sigma_{i}^{2}}}.
\]
Since this does not depend on \(\mathbf{X}\), we can pull it out of the expectation:
\[
\mathbb{E}\big[(\Theta - \hat{\Theta}(\mathbf{X}))^{2}\big]
=
\frac{1}{\displaystyle \frac{1}{\sigma_{0}^{2}} + \sum_{i=1}^{n}\frac{1}{\sigma_{i}^{2}}}
\cdot
\mathbb{E} [statswithr.github](https://statswithr.github.io/book/inference-and-decision-making-with-multiple-parameters.html)
=
\frac{1}{\displaystyle \frac{1}{\sigma_{0}^{2}} + \sum_{i=1}^{n}\frac{1}{\sigma_{i}^{2}}}.
\]
Thus the unconditional mean squared error equals the conditional mean squared error.  
In other words, the expected performance of the estimator, before making any measurements, is equal to the posterior variance obtained after the measurements, and this value is the same for every possible outcome of the measurements.

Let us now interpret this expression and examine some special cases.  
First, suppose that at least one of the noise variances \(\sigma_{i}^{2}\) is very small.  
Then the corresponding term \(1/\sigma_{i}^{2}\) is very large.  
Consequently, the sum
\[
\frac{1}{\sigma_{0}^{2}} + \sum_{i=1}^{n}\frac{1}{\sigma_{i}^{2}}
\]
is large, and its reciprocal, the mean squared error, is small.  
This reflects the fact that if we have even a single high-quality (low-noise) measurement, we can estimate \(\Theta\) quite accurately.

On the other hand, if all noise variances \(\sigma_{i}^{2}\) are large, then each term \(1/\sigma_{i}^{2}\) is small.  
The sum of these small terms is also small, so its reciprocal is large.  
In this case, all measurements are very noisy, and we cannot expect to estimate \(\Theta\) well; the mean squared error is large.

Next, consider the special case where all variances, including the prior variance, are equal:
\[
\sigma_{0}^{2}
=
\sigma_{1}^{2}
=
\dots
=
\sigma_{n}^{2}
=
\sigma^{2}.
\]
Then
\[
\frac{1}{\sigma_{0}^{2}} + \sum_{i=1}^{n}\frac{1}{\sigma_{i}^{2}}
=
\frac{1}{\sigma^{2}} + \sum_{i=1}^{n}\frac{1}{\sigma^{2}}
=
\frac{n+1}{\sigma^{2}}.
\]
Therefore, the mean squared error simplifies to
\[
\mathbb{E}\big[(\Theta - \hat{\Theta}(\mathbf{X}))^{2}\big]
=
\frac{1}{(n+1)/\sigma^{2}}
=
\frac{\sigma^{2}}{n+1}.
\]
This expression shows explicitly how the estimation error decreases as the number of observations \(n\) increases.  
Each additional measurement contributes an extra term \(1/\sigma^{2}\) to the denominator, reducing the overall variance and improving the quality of the estimate.

\bigskip
\textsc{A concrete example: single observation and standard normals}

To better appreciate the invariance of the conditional variance with respect to the data, let us revisit the simple model with a single observation and standard normals.  
Assume
\[
\Theta \sim \mathcal{N}(0,1),
\qquad
W \sim \mathcal{N}(0,1),
\qquad
\Theta \perp W,
\]
and
\[
X
=
\Theta + W.
\]
In this case, \(n=1\), \(\sigma_{0}^{2} = 1\), and \(\sigma_{1}^{2} = 1\).  
The general formula gives
\[
\operatorname{Var}(\Theta \mid X = x)
=
\frac{1}{\displaystyle \frac{1}{1} + \frac{1}{1}}
=
\frac{1}{2}.
\]
Thus, starting from a prior variance of \(1\), the variance after observing \(X\) is reduced to \(1/2\).  
This reduction in variance reflects the information gained from the measurement.

Importantly, the conditional variance is equal to \(1/2\) for every value of \(x\).  
Whether the observed value \(x\) is close to \(0\) or far away from \(0\), the remaining uncertainty about \(\Theta\), as measured by the posterior variance, is the same.  
If we observe \(x=0\), our estimate is \(\hat{\theta} = x/2 = 0\), and the posterior distribution of \(\Theta\) is \(\mathcal{N}(0,1/2)\), which is narrower than the prior \(\mathcal{N}(0,1)\).  
If instead we observe, say, \(x\) far from \(0\), then the estimate \(\hat{\theta} = x/2\) will be correspondingly far from \(0\), and the posterior distribution will be \(\mathcal{N}(x/2,1/2)\).  
In this latter case, the observation is surprising in light of the prior, but once it is observed, the remaining uncertainty about \(\Theta\) is again characterized by a normal distribution with variance \(1/2\).

This behavior is a distinctive feature of linear Gaussian models.  
The combination of normal priors with normal likelihoods leads to posterior distributions that are also normal, with a variance that depends only on the prior and noise variances, and not on the particular observed values.  
As a result, the performance of the MAP or least mean squares estimator can be fully characterized by a single number, the posterior variance
\[
\frac{1}{\displaystyle \frac{1}{\sigma_{0}^{2}} + \sum_{i=1}^{n}\frac{1}{\sigma_{i}^{2}}},
\]
which serves both as the conditional and unconditional mean squared error.



\subsubsection{Multiple parameters; trajectory estimation}

We now consider a model with multiple parameters and multiple observations, motivated by a simple physical scenario.  
A ball is thrown upwards and follows a trajectory determined by Newton's laws.  
The vertical displacement \(x(t)\) of the ball at time \(t\) is modeled by a quadratic function
\[
x(t)
=
\Theta_{0} + \Theta_{1} t + \Theta_{2} t^{2},
\]
where:
\[
\Theta_{0} \text{ is the initial height},
\quad
\Theta_{1} \text{ is the initial vertical velocity},
\quad
\Theta_{2} \text{ is related to the gravitational acceleration (typically negative)}.
\]

Suppose that the parameters \(\Theta_{0},\Theta_{1},\Theta_{2}\) are unknown.  
We model them as random variables and treat the actual physical values as realizations of these random variables.  
To apply Bayesian inference, we specify a joint prior distribution for the parameter vector
\[
\boldsymbol{\Theta}
=
(\Theta_{0},\Theta_{1},\Theta_{2})^{\top}.
\]
A common assumption is that the parameters are independent and normally distributed:
\[
\Theta_{k} \sim \mathcal{N}(0,\sigma_{\Theta_{k}}^{2}),
\quad
k=0,1,2,
\]
with independence across \(k\).  
The prior variances \(\sigma_{\Theta_{k}}^{2}\) are chosen to reflect our initial uncertainty about the parameters (for example, how well we know the initial height or velocity).

We observe the ball at times \(t_{1},\dots,t_{n}\).  
At each time \(t_{i}\), we measure the vertical position with noise, obtaining
\[
X_{i}
=
\Theta_{0} + \Theta_{1} t_{i} + \Theta_{2} t_{i}^{2} + W_{i},
\quad
i = 1,\dots,n,
\]
where \(W_{i}\) is the measurement noise.  
We assume that
\[
W_{i} \sim \mathcal{N}(0,\sigma_{i}^{2}),
\]
and that the noise variables \(W_{1},\dots,W_{n}\) are independent of each other and independent of \(\boldsymbol{\Theta}\).  
This is often a realistic model for repeated measurements made by a sensor.

Let
\[
\mathbf{X}
=
(X_{1},\dots,X_{n})^{\top},
\quad
\mathbf{x}
=
(x_{1},\dots,x_{n})^{\top},
\]
and
\[
\boldsymbol{\Theta}
=
(\Theta_{0},\Theta_{1},\Theta_{2})^{\top},
\quad
\boldsymbol{\theta}
=
(\theta_{0},\theta_{1},\theta_{2})^{\top}.
\]
Our goal is to estimate \(\boldsymbol{\Theta}\) from the observations \(\mathbf{X}\) using Bayesian inference, that is, to compute the posterior density \(f_{\boldsymbol{\Theta}\mid \mathbf{X}}(\boldsymbol{\theta}\mid \mathbf{x})\).  
Bayes' rule for continuous random vectors gives
\[
f_{\boldsymbol{\Theta}\mid \mathbf{X}}(\boldsymbol{\theta}\mid \mathbf{x})
=
\frac{f_{\mathbf{X}\mid \boldsymbol{\Theta}}(\mathbf{x}\mid \boldsymbol{\theta})\,f_{\boldsymbol{\Theta}}(\boldsymbol{\theta})}{f_{\mathbf{X}}(\mathbf{x})}.
\]
Here:
\begin{itemize}
\item \(f_{\boldsymbol{\Theta}}(\boldsymbol{\theta})\) is the prior density of the parameter vector,
\item \(f_{\mathbf{X}\mid \boldsymbol{\Theta}}(\mathbf{x}\mid \boldsymbol{\theta})\) is the conditional joint density of the measurements given the parameters,
\item \(f_{\mathbf{X}}(\mathbf{x})\) is the marginal density of the observations.
\end{itemize}

We first determine the conditional distribution of a single observation \(X_{i}\) given \(\boldsymbol{\Theta} = \boldsymbol{\theta}\).  
Conditioning on \(\boldsymbol{\Theta} = \boldsymbol{\theta}\), the measurement model becomes
\[
X_{i} \mid (\boldsymbol{\Theta} = \boldsymbol{\theta})
=
\theta_{0} + \theta_{1} t_{i} + \theta_{2} t_{i}^{2} + W_{i}.
\]
The noise term \(W_{i}\) remains \(\mathcal{N}(0,\sigma_{i}^{2})\) under this conditioning because \(W_{i}\) is independent of \(\boldsymbol{\Theta}\).  
Therefore,
\[
X_{i} \mid (\boldsymbol{\Theta} = \boldsymbol{\theta})
\sim
\mathcal{N}\big(\theta_{0} + \theta_{1} t_{i} + \theta_{2} t_{i}^{2},\ \sigma_{i}^{2}\big),
\]
and the conditional density is
\[
f_{X_{i}\mid \boldsymbol{\Theta}}(x_{i}\mid \boldsymbol{\theta})
=
\frac{1}{\sqrt{2\pi\sigma_{i}^{2}}}
\exp\Bigg(
-\frac{\big(x_{i} - (\theta_{0} + \theta_{1} t_{i} + \theta_{2} t_{i}^{2})\big)^{2}}{2\sigma_{i}^{2}}
\Bigg).
\]

Next, we determine the conditional joint density \(f_{\mathbf{X}\mid \boldsymbol{\Theta}}(\mathbf{x}\mid \boldsymbol{\theta})\).  
Given \(\boldsymbol{\Theta} = \boldsymbol{\theta}\), each \(X_{i}\) is obtained from the corresponding noise variable \(W_{i}\) by adding the deterministic quantity \(\theta_{0} + \theta_{1} t_{i} + \theta_{2} t_{i}^{2}\).  
The noise variables \(W_{1},\dots,W_{n}\) are independent, so the measurements \(X_{1},\dots,X_{n}\) are also independent in the conditional universe where \(\boldsymbol{\Theta}\) is fixed.  
It follows that
\[
f_{\mathbf{X}\mid \boldsymbol{\Theta}}(\mathbf{x}\mid \boldsymbol{\theta})
=
\prod_{i=1}^{n}
f_{X_{i}\mid \boldsymbol{\Theta}}(x_{i}\mid \boldsymbol{\theta})
=
\prod_{i=1}^{n}
\frac{1}{\sqrt{2\pi\sigma_{i}^{2}}}
\exp\Bigg(
-\frac{\big(x_{i} - (\theta_{0} + \theta_{1} t_{i} + \theta_{2} t_{i}^{2})\big)^{2}}{2\sigma_{i}^{2}}
\Bigg).
\]

For the prior, assuming independent normal components with zero means, we have
\[
f_{\boldsymbol{\Theta}}(\boldsymbol{\theta})
=
\prod_{k=0}^{2}
\frac{1}{\sqrt{2\pi\sigma_{\Theta_{k}}^{2}}}
\exp\Big(-\frac{\theta_{k}^{2}}{2\sigma_{\Theta_{k}}^{2}}\Big).
\]
Putting these pieces into Bayes' rule and ignoring the marginal density \(f_{\mathbf{X}}(\mathbf{x})\), which does not depend on \(\boldsymbol{\theta}\), we obtain, up to a normalizing constant,
\[
f_{\boldsymbol{\Theta}\mid \mathbf{X}}(\boldsymbol{\theta}\mid \mathbf{x})
\propto
f_{\mathbf{X}\mid \boldsymbol{\Theta}}(\mathbf{x}\mid \boldsymbol{\theta})\,f_{\boldsymbol{\Theta}}(\boldsymbol{\theta})
\]
\[
\propto
\Bigg[
\prod_{i=1}^{n}
\frac{1}{\sqrt{2\pi\sigma_{i}^{2}}}
\exp\Bigg(
-\frac{\big(x_{i} - (\theta_{0} + \theta_{1} t_{i} + \theta_{2} t_{i}^{2})\big)^{2}}{2\sigma_{i}^{2}}
\Bigg)
\Bigg]
\Bigg[
\prod_{k=0}^{2}
\frac{1}{\sqrt{2\pi\sigma_{\Theta_{k}}^{2}}}
\exp\Big(-\frac{\theta_{k}^{2}}{2\sigma_{\Theta_{k}}^{2}}\Big)
\Bigg].
\]
The multiplicative constants in front of the exponentials do not depend on \(\boldsymbol{\theta}\) and can be absorbed into a single constant.  
Therefore, the posterior density is proportional to
\[
\exp\Bigg(
-\sum_{k=0}^{2}\frac{\theta_{k}^{2}}{2\sigma_{\Theta_{k}}^{2}}
-
\sum_{i=1}^{n}
\frac{\big(x_{i} - (\theta_{0} + \theta_{1} t_{i} + \theta_{2} t_{i}^{2})\big)^{2}}{2\sigma_{i}^{2}}
\Bigg).
\]

This is the exponential of the negative of a quadratic function of \(\boldsymbol{\theta}\).  
Because the exponent is a quadratic form in the parameters, the posterior distribution of \(\boldsymbol{\Theta}\) is a multivariate normal distribution.  
Although we do not explicitly write its mean vector and covariance matrix here, the structure matches the general result: a normal prior combined with a linear-Gaussian observation model produces a normal posterior.

To obtain a point estimate of the parameters, a natural choice is the maximum a posteriori (MAP) estimate.  
The MAP estimate
\[
\hat{\boldsymbol{\theta}}_{\text{MAP}}
=
\arg\max_{\boldsymbol{\theta}} f_{\boldsymbol{\Theta}\mid \mathbf{X}}(\boldsymbol{\theta}\mid \mathbf{x})
\]
is the value of \(\boldsymbol{\theta}\) that maximizes the posterior density.  
Since the exponential function is strictly increasing, maximizing the posterior is equivalent to minimizing the negative of the exponent.  
Thus we consider the quadratic cost function
\[
J(\boldsymbol{\theta})
=
\sum_{k=0}^{2}\frac{\theta_{k}^{2}}{2\sigma_{\Theta_{k}}^{2}}
+
\sum_{i=1}^{n}
\frac{\big(x_{i} - (\theta_{0} + \theta_{1} t_{i} + \theta_{2} t_{i}^{2})\big)^{2}}{2\sigma_{i}^{2}}.
\]
The MAP estimate is
\[
\hat{\boldsymbol{\theta}}_{\text{MAP}}
=
\arg\min_{\boldsymbol{\theta}} J(\boldsymbol{\theta}).
\]

The function \(J(\boldsymbol{\theta})\) is a quadratic function of \(\theta_{0},\theta_{1},\theta_{2}\).  
To find its minimum, we take the partial derivatives with respect to each parameter and set them equal to zero:
\[
\frac{\partial J}{\partial \theta_{k}}(\boldsymbol{\theta}) = 0,
\quad
k = 0,1,2.
\]
Each partial derivative is a linear function of \(\theta_{0},\theta_{1},\theta_{2}\), because the derivative of a quadratic function is linear.  
Therefore, we obtain a system of three linear equations in the three unknowns \(\theta_{0},\theta_{1},\theta_{2}\).  
In matrix form, this can be written as
\[
\mathbf{A}\boldsymbol{\theta}
=
\mathbf{b},
\]
for an appropriate \(3 \times 3\) matrix \(\mathbf{A}\) and vector \(\mathbf{b}\) determined by the times \(t_{i}\), the observations \(x_{i}\), the noise variances \(\sigma_{i}^{2}\), and the prior variances \(\sigma_{\Theta_{k}}^{2}\).  
Solving this linear system yields the MAP estimate \(\hat{\boldsymbol{\theta}}_{\text{MAP}}\).

In practice, this solution is computed numerically using standard linear algebra methods.  
The resulting estimator has several important properties that extend those seen in the single-parameter case.  
The posterior distribution remains Gaussian, the MAP estimate coincides with the posterior mean, and the estimator is a linear function of the observations \(\mathbf{X}\).  
These features make such linear-Gaussian models particularly attractive for applications such as trajectory estimation and tracking.


